\subsection{Affine Maps and Affine Groups}

Throughout, $\mathcal{V}$ is finite-dimensional over a fixed field $\mathbb{F}$. For $n\in\mathbb{Z}_{>0}$, put $\mathcal{C}_n\defeq\{\mathbf{r}=(r_1,\ldots,r_n)\in\mathbb{F}^n:\sum\limits_{j=1}^n r_j=1\}$.

\begin{definition} \label{def:affine-transformation} \hlt{Affine Transformation}\\
A map $f:\mathcal{V}\to\mathcal{V}$ is \hlt{affine} iff $(\forall\,n\in\mathbb{Z}_{>0})(\forall\,(\mathbf{x}_1,\ldots,\mathbf{x}_n)\in\mathcal{V}^n)(\forall\,\mathbf{r}\in\mathcal{C}_n)\,f(\sum\limits_{j=1}^n r_j\mathbf{x}_j)=\sum\limits_{j=1}^n r_jf(\mathbf{x}_j)$.
\end{definition}

\begin{theorem} \label{thm:affine-map-two-point-test} \hlt{Two-Point Test for Affine Maps}\\
Assume $\operatorname{char}(\mathbb{F})\neq2$. A map $f:\mathcal{V}\to\mathcal{V}$ is affine iff $(\forall\,\mathbf{x},\mathbf{y}\in\mathcal{V})(\forall\,r\in\mathbb{F})\,f(r\mathbf{x}+(1-r)\mathbf{y})=rf(\mathbf{x})+(1-r)f(\mathbf{y})$.
\begin{proof}
\begin{enumerate}[label=\roman*.]
\setlength{\itemsep}{0pt}
\item[$(\Rightarrow)$] The identity holds by Definition \ref{def:affine-transformation} with $n=2$ and $(r,1-r)\in\mathcal{C}_2$.
\item[$(\Leftarrow)$] Put $\Gamma_f\defeq\{(\mathbf{x},f(\mathbf{x})):\mathbf{x}\in\mathcal{V}\}\subseteq\mathcal{V}\times\mathcal{V}$.\\
The relation $(\forall\,\mathbf{x},\mathbf{y}\in\mathcal{V})(\forall\,r\in\mathbb{F})\,r(\mathbf{x},f(\mathbf{x}))+(1-r)(\mathbf{y},f(\mathbf{y}))\in\Gamma_f$ holds by the hypothesis, so $\Gamma_f$ is two-affine closed.\\
The nonempty set $\Gamma_f$ is affine closed by Theorem \ref{thm:affine-closed-two-affine-closed} in $\mathcal{V}\times\mathcal{V}$.\\
Hence $(\forall\,n\in\mathbb{Z}_{>0})(\forall\,(\mathbf{x}_1,\ldots,\mathbf{x}_n)\in\mathcal{V}^n)(\forall\,\mathbf{r}\in\mathcal{C}_n)\,(\sum\limits_{j=1}^n r_j\mathbf{x}_j,\sum\limits_{j=1}^n r_jf(\mathbf{x}_j))\in\Gamma_f$.\\
By the definition of $\Gamma_f$, this membership is equivalent to the identity in Definition \ref{def:affine-transformation}; hence $f$ is affine.
\end{enumerate}
\hlt{Structural Commentary}: Transfer binary preservation to the graph, then promote its two-affine closure to affine closure.\\
\hlt{Key Insight}: Binary preservation is two-affine closure of the graph.
\end{proof}
\end{theorem}

\begin{definition} \label{def:translation-map} \hlt{Translations}\\
For $\mathbf{a}\in\mathcal{V}$, the \hlt{translation by $\mathbf{a}$} is the map $T_{\mathbf{a}}:\mathcal{V}\to\mathcal{V}$ defined by $T_{\mathbf{a}}(\mathbf{x})=\mathbf{x}+\mathbf{a}$.
\end{definition}

\begin{theorem} \label{thm:affine-map-linear-part-translation} \hlt{Affine Maps Are Linear Maps Followed by Translations}\\
A map $f:\mathcal{V}\to\mathcal{V}$ is affine iff $(\exists!\,(\mathbf{a},\tau)\in\mathcal{V}\times\LL(\mathcal{V}))\,[\mathbf{a}=f(\mathbf{0})\land\tau=T_{-\mathbf{a}}\circ f\land f=T_{\mathbf{a}}\circ\tau]$.
\begin{proof}
\begin{enumerate}[label=\roman*.]
\setlength{\itemsep}{0pt}
\item[$(\Rightarrow)$] Put $\mathbf{a}\defeq f(\mathbf{0})$ and $\tau\defeq T_{-\mathbf{a}}\circ f$.\\
For $r,s\in\mathbb{F}$ and $\mathbf{x},\mathbf{y}\in\mathcal{V}$, $\tau(r\mathbf{x}+s\mathbf{y})=r\tau(\mathbf{x})+s\tau(\mathbf{y})$ by Definition \ref{def:affine-transformation} with $(\mathbf{x},\mathbf{y},\mathbf{0})$ and $(r,s,1-r-s)\in\mathcal{C}_3$.\\
Thus $\tau\in\LL(\mathcal{V})$, $\mathbf{a}=f(\mathbf{0})$, $\tau=T_{-\mathbf{a}}\circ f$, and $f=T_{\mathbf{a}}\circ\tau$.\\
If $f=T_{\mathbf{b}}\circ\sigma$ with $(\mathbf{b},\sigma)\in\mathcal{V}\times\LL(\mathcal{V})$, then $\mathbf{b}=f(\mathbf{0})=\mathbf{a}$ and $\sigma=T_{-\mathbf{b}}\circ f=\tau$, so the pair is unique.
\item[$(\Leftarrow)$] Suppose $(\mathbf{a},\tau)\in\mathcal{V}\times\LL(\mathcal{V})$ satisfies $f=T_{\mathbf{a}}\circ\tau$.\\
For $n\in\mathbb{Z}_{>0}$, $(\mathbf{x}_1,\ldots,\mathbf{x}_n)\in\mathcal{V}^n$, and $\mathbf{r}\in\mathcal{C}_n$, $f(\sum\limits_{j=1}^n r_j\mathbf{x}_j)=\mathbf{a}+\sum\limits_{j=1}^n r_j\tau(\mathbf{x}_j)=\sum\limits_{j=1}^n r_jf(\mathbf{x}_j)$. Therefore $f$ affine by Definition \ref{def:affine-transformation}.
\end{enumerate}
\hlt{Structural Commentary}: Translate $f(\mathbf{0})$ to zero, verify linearity, and identify both factors by evaluation at zero.\\
\hlt{Key Insight}: The value $f(\mathbf{0})$ determines the translation factor.
\end{proof}
\end{theorem}

\begin{definition} \label{def:affine-group-translations-linear-part} \hlt{Affine Group Data}
\begin{enumerate}[label=(\roman*)]
\setlength{\itemsep}{0pt}
\item The \hlt{general linear group} is $\operatorname{GL}(\mathcal{V})\defeq\{\tau\in\LL(\mathcal{V}):(\exists\,\sigma\in\LL(\mathcal{V}))\,[\sigma\tau=I_{\mathcal{V}}=\tau\sigma]\}$.
\item The \hlt{affine group} is $\operatorname{Aff}(\mathcal{V})\defeq\{T_{\mathbf{a}}\circ\tau:\mathbf{a}\in\mathcal{V},\,\tau\in\operatorname{GL}(\mathcal{V})\}$.
\item The \hlt{translation group} is $\operatorname{Trans}(\mathcal{V})\defeq\{T_{\mathbf{a}}:\mathbf{a}\in\mathcal{V}\}$.
\item The \hlt{linear-part map} is $\Lambda:\operatorname{Aff}(\mathcal{V})\to\operatorname{GL}(\mathcal{V})$, defined by $\Lambda(T_{\mathbf{a}}\circ\tau)=\tau$.
\end{enumerate}
Well-definedness holds by Theorem \ref{thm:affine-map-linear-part-translation}.
\end{definition}

\begin{corollary} \label{cor:affine-map-bijective-linear-part} \hlt{Bijective Affine Maps}\\
Let $\mathbf{a}\in\mathcal{V}$ and $\tau\in\LL(\mathcal{V})$. Then $T_{\mathbf{a}}\circ\tau$ is bijective iff $\tau\in\operatorname{GL}(\mathcal{V})$.
\begin{proof}
\begin{enumerate}[label=\roman*.]
\setlength{\itemsep}{0pt}
\item[$(\Rightarrow)$] The map $T_{-\mathbf{a}}$ is bijective with inverse $T_{\mathbf{a}}$. Thus $\tau=T_{-\mathbf{a}}\circ(T_{\mathbf{a}}\circ\tau)$ is bijective, so $\tau\in\operatorname{GL}(\mathcal{V})$ by Proposition \ref{prop:linear-map-invertible-iff-bijective}.
\item[$(\Leftarrow)$] If $\tau\in\operatorname{GL}(\mathcal{V})$, then $\tau^{-1}\circ T_{-\mathbf{a}}$ is a two-sided inverse of $T_{\mathbf{a}}\circ\tau$. Hence $T_{\mathbf{a}}\circ\tau$ is bijective by Proposition \ref{prop:linear-map-invertible-iff-bijective}.
\end{enumerate}
\hlt{Structural Commentary}: Cancel the translation forward and compose the linear and translation inverses backward.\\
\hlt{Key Insight}: Translation does not affect bijectivity.
\end{proof}
\end{corollary}

\begin{theorem} \label{thm:affine-group-linear-part} \hlt{Affine Group and Translations}
\begin{enumerate}[label=(\roman*)]
\setlength{\itemsep}{0pt}
\item $(\forall\,\mathbf{a},\mathbf{b}\in\mathcal{V})(\forall\,\tau,\sigma\in\LL(\mathcal{V}))\,(T_{\mathbf{a}}\circ\tau)\circ(T_{\mathbf{b}}\circ\sigma)=T_{\mathbf{a}+\tau(\mathbf{b})}\circ(\tau\sigma)$.
\item $\operatorname{GL}(\mathcal{V})$ is a group under composition.
\item $\operatorname{Aff}(\mathcal{V})$ is a group under composition, with identity $T_{\mathbf{0}}\circ I_{\mathcal{V}}$ and $(T_{\mathbf{a}}\circ\tau)^{-1}=T_{-\tau^{-1}(\mathbf{a})}\circ\tau^{-1}$.
\item $\operatorname{Trans}(\mathcal{V})\trianglelefteq\operatorname{Aff}(\mathcal{V})$.
\item $\Lambda$ is a group homomorphism.
\item $\Lambda$ is surjective.
\item $\ker\Lambda=\operatorname{Trans}(\mathcal{V})$.
\item $\operatorname{Aff}(\mathcal{V})/\operatorname{Trans}(\mathcal{V})\cong\operatorname{GL}(\mathcal{V})$.
\end{enumerate}
\begin{proof}
\begin{enumerate}[label=(\roman*)]
\setlength{\itemsep}{0pt}
\item For $\mathbf{x}\in\mathcal{V}$, $((T_{\mathbf{a}}\circ\tau)\circ(T_{\mathbf{b}}\circ\sigma))(\mathbf{x})=\mathbf{a}+\tau(\mathbf{b})+(\tau\sigma)(\mathbf{x})$. Hence $(T_{\mathbf{a}}\circ\tau)\circ(T_{\mathbf{b}}\circ\sigma)=T_{\mathbf{a}+\tau(\mathbf{b})}\circ(\tau\sigma)$ by Assumption \ref{assump:function-extensionality}.
\item The identity $I_{\mathcal{V}}$ belongs to $\operatorname{GL}(\mathcal{V})$, and composition is associative. Closure holds by Proposition \ref{prop:inverse-of-product-linear-maps}(i).\\
If $\tau\in\operatorname{GL}(\mathcal{V})$, then $\tau^{-1}\in\operatorname{GL}(\mathcal{V})$ by Definition \ref{def:invertible-linear-map-and-isomorphism} and Proposition \ref{prop:inverse-linear-map-is-unique}. Thus $\operatorname{GL}(\mathcal{V})$ is a group under composition.
\item Closure holds by (i)--(ii), associativity is inherited from composition, and the identity is $T_{\mathbf{0}}\circ I_{\mathcal{V}}$.\\
For $g=T_{\mathbf{a}}\circ\tau$, put $h\defeq T_{-\tau^{-1}(\mathbf{a})}\circ\tau^{-1}\in\operatorname{Aff}(\mathcal{V})$. The identity $gh=hg=T_{\mathbf{0}}\circ I_{\mathcal{V}}$ holds by (i), so $h=g^{-1}$ and $\operatorname{Aff}(\mathcal{V})$ is a group.
\item The identities $T_{\mathbf{a}}T_{\mathbf{b}}=T_{\mathbf{a}+\mathbf{b}}$, $T_{\mathbf{0}}=I_{\mathcal{V}}$, and $T_{\mathbf{a}}^{-1}=T_{-\mathbf{a}}$ hold by (i), so $\operatorname{Trans}(\mathcal{V})\leq\operatorname{Aff}(\mathcal{V})$.\\
If $g=T_{\mathbf{a}}\circ\tau$, then $(\forall\,\mathbf{b}\in\mathcal{V})\,gT_{\mathbf{b}}g^{-1}=T_{\tau(\mathbf{b})}$ by (i). The map $\tau$ is bijective by Proposition \ref{prop:linear-map-invertible-iff-bijective}, so $g\operatorname{Trans}(\mathcal{V})g^{-1}=\operatorname{Trans}(\mathcal{V})$.
\item For $g=T_{\mathbf{a}}\circ\tau$ and $h=T_{\mathbf{b}}\circ\sigma$, $\Lambda(gh)=\tau\sigma=\Lambda(g)\Lambda(h)$ by (i).
\item $(\forall\,\tau\in\operatorname{GL}(\mathcal{V}))\,\Lambda(T_{\mathbf{0}}\circ\tau)=\tau$, so $\Lambda$ is surjective.
\item The identity $\ker\Lambda=\{T_{\mathbf{a}}\circ I_{\mathcal{V}}:\mathbf{a}\in\mathcal{V}\}=\operatorname{Trans}(\mathcal{V})$ holds by Definition \ref{def:affine-group-translations-linear-part}.
\item Put $N\defeq\operatorname{Trans}(\mathcal{V})$ and define $\overline{\Lambda}(gN)\defeq\Lambda(g)$. If $gN=hN$, then $h^{-1}g\in N=\ker\Lambda$.\\
The identities $\Lambda(h^{-1})=\Lambda(h)^{-1}$ and $\Lambda(h)^{-1}\Lambda(g)=I_{\mathcal{V}}$ hold by (v), so $\Lambda(g)=\Lambda(h)$ and $\overline{\Lambda}$ is well defined.\\
For $g,h\in\operatorname{Aff}(\mathcal{V})$, $\overline{\Lambda}((gN)(hN))=\Lambda(gh)=\Lambda(g)\Lambda(h)=\overline{\Lambda}(gN)\overline{\Lambda}(hN)$ by (v), so $\overline{\Lambda}$ is a homomorphism.\\
It is surjective by (vi). If $\overline{\Lambda}(gN)=\overline{\Lambda}(hN)$, then $\Lambda(h^{-1}g)=I_{\mathcal{V}}$ by (v), so $h^{-1}g\in N$ by (vii) and $gN=hN$.\\
Hence $\overline{\Lambda}$ is a bijective homomorphism and $\operatorname{Aff}(\mathcal{V})/N\cong\operatorname{GL}(\mathcal{V})$.
\end{enumerate}
\hlt{Structural Commentary}: Product and inverse laws establish normal translations before $\Lambda$ descends injectively to their quotient.\\
\hlt{Key Insight}: The kernel of the linear-part map is exactly translation.
\end{proof}
\end{theorem}
